% R12 proofs of the whitened-innovation laws. Numerical checks: research/r7c/r12/verify_laws_output.txt.
\section{Proofs of the whitened-innovation laws}\label{sup:r12proofs}
Throughout, $z\sim\mathcal{CN}(0,I)$ denotes whitened clutter innovations (unit variance after dividing by $\sigma_e$), $g\sim\mathcal{CN}(0,1)$ the Swerling-1 amplitude, and $u$ the whitened target signature, so that the whitened episode is $z+gu\sim\mathcal{CN}(0,I+uu^{\mathsf H})$. $E$, $E_\pm$ are standard exponentials and $\Gamma_k$ a Gamma$(k,1)$ variable, all independent.

\subsection{Proposition (post-onset rank-one law)}
Write the detection event as $F>0$ with $F=x^{\mathsf H}Mx$, $x=z+gu$ and $M=\operatorname{diag}(-\beta I_m,1)$, $\beta=q^2/m$. As in Appendix~\ref{M-app:law} of the paper, $F\overset{d}{=}\sum_k\lambda_kE_k$ with $\lambda_k$ the eigenvalues of $C^{1/2}MC^{1/2}$, $C=I+uu^{\mathsf H}$, which are those of $MC=M+(Mu)u^{\mathsf H}$. For a rank-one update of a diagonal matrix,
\[
\det(\lambda I-M-Muu^{\mathsf H})=\det(\lambda I-M)\Big(1-\sum_i\frac{M_{ii}|u_i|^2}{\lambda-M_{ii}}\Big)=(\lambda-1)(\lambda+\beta)^m\Big(1-\frac{\mathcal E_0}{\lambda-1}+\frac{\beta \mathcal E_H}{\lambda+\beta}\Big),
\]
with $\mathcal E_0=|u_Y|^2$ and $\mathcal E_H=\|u_H\|^2$. Hence $-\beta$ is an eigenvalue of multiplicity $m-1$ and the other two solve $(\lambda-1)(\lambda+\beta)-\mathcal E_0(\lambda+\beta)+\beta\mathcal E_H(\lambda-1)=0$, which is the quadratic of the proposition. Their product is $-\beta(1+\mathcal E_0+\mathcal E_H)<0$, so $\lambda_+>0>\lambda_-$. Then
\[
\PP\{F>0\}=\PP\{\lambda_+E_+>|\lambda_-|E_-+\beta\Gamma_{m-1}\}=\mathbb E\,e^{-(|\lambda_-|E_-+\beta\Gamma_{m-1})/\lambda_+}=\frac{\lambda_+}{\lambda_++|\lambda_-|}\Big(1+\frac{\beta}{\lambda_+}\Big)^{-(m-1)}.
\]
For $\mathcal E_H=0$ the roots are $1+\mathcal E_0$ and $-\beta$, and the expression reduces to $(1+\beta/(1+\mathcal E_0))^{-m}$. The signature of a persistent target follows from the whitened amplitudes in \eqref{M-eq:whitenedscr}: the onset innovation carries $S_w$, every later one $S_w|d|^2$.

\subsection{Corollary (visibility horizon)}
With $\mathcal E_0=S_wD$ and $\mathcal E_H=S_w(1+(\ell-1)D)$, $D=|d(\omega)|^2$, the linear coefficient is $B=1-\beta+S_w\chi$, $\chi=D-\beta\{1+(\ell-1)D\}$, and $\lambda_\pm=\{B\pm(B^2+4\beta(1+\mathcal E_0+\mathcal E_H))^{1/2}\}/2$ with $1+\mathcal E_0+\mathcal E_H=O(S_w)$. If $\chi>0$, then $\lambda_+\sim S_w\chi\to\infty$ while $\lambda_-$ stays bounded, so $P_{\rm d}\to1$. If $\chi<0$, then $\lambda_+\to\beta(1+\ell D)/|\chi|$ stays bounded while $\lambda_-\sim S_w\chi\to-\infty$, so $\lambda_+/(\lambda_+-\lambda_-)\to0$ and $P_{\rm d}\to0$. The sign change $\chi=0$ is $\ell=1+1/\beta-1/D$. At equality, $\lambda_\pm\sim\pm\sqrt{\beta S_w(1+\ell D)}$ and $P_{\rm d}\to1/2$. The horizon describes the post-onset interval in which the onset innovation remains in the history, not look zero. With a guard, only the history energy is delayed. Within protected looks $1,\ldots,\Delta$, $P_{\rm d}\to1$ as $S\to\infty$ for $D>0$. The critical value $\Delta+\ell^*$ applies beyond those looks; if $\ell^*<1$, loss first occurs at $\Delta+1$ (at $\ell^*=1$ that look tends to $1/2$).

\subsection{Order-statistic scale}
Conditional on $G=|g|^2$, the history powers $|z_j+gu_j|^2$ are independent, exponential for clean innovations and non-central (with $2|z_j+gu_j|^2\sim\chi^2_2(2G|u_j|^2)$) for contaminated ones. The distribution function of their $k$-th smallest value $X_{(k)}$ is the Poisson-binomial probability that at least $k$ are below $x$, and $P_{\rm d}=\mathbb E_G\int\PP\{|z_Y+gu_Y|^2>q^2x\mid G\}\,dF_{X_{(k)}\mid G}(x)$. The code evaluates this by midpoint quadrature in the distribution function of $G$ and a dynamic program for the Poisson-binomial law. If $\ell\le m-k$ contaminated powers diverge, $X_{(k)}$ is the $k$-th smallest of the $m-\ell$ clean ones, which is independent of $z_Y$, and the OS-CFAR law with $m-\ell$ cells gives \eqref{M-eq:osstrong}. If $\ell>m-k$, let $j=k-(m-\ell)$ and let $a_{(j)}$ be the $j$-th smallest of $\{1,D,\ldots,D\}$ ($\ell$ entries), $D=|d|^2$. The tested-power/scale ratio converges almost surely to $D/a_{(j)}$. Hence $P_{\rm d}\to0$ when $q^2>D/a_{(j)}$, in particular when $q^2>\max\{1,D\}$, as for the reported $k=8,12$ settings at $\alpha=0.01$. Outside this condition self-masking is not assured.

\subsection{Dwell integration}
Under $H_0$ the event is $\Gamma_K>\tau\Gamma_m$, and conditioning on $\Gamma_m$ with the Poisson form of the Gamma tail gives $\sum_{i<K}\mathbb E[e^{-\tau\Gamma_m}(\tau\Gamma_m)^i/i!]=\sum_{i<K}\binom{m+i-1}{i}\tau^i(1+\tau)^{-(m+i)}$. With a Swerling-1 target the dwell form has one eigenvalue $\delta=1+\mathcal E_D$ and $K-1$ unit ones. For $v>0$,
\[
\PP\{\delta E+\Gamma_{K-1}>v\}=\PP\{\Gamma_{K-1}>v\}+e^{-v/\delta}(1-1/\delta)^{-(K-1)}\PP\{\Gamma_{K-1}\le v(1-1/\delta)\}.
\]
Averaging over $v=\tau\Gamma_m$ uses $\mathbb E e^{-s\Gamma_m}=(1+s)^{-m}$ and $\mathbb E[e^{-s\Gamma_m}\PP\{\Gamma_{K-1}>\tau'\Gamma_m\mid\Gamma_m\}]=\sum_{i\le K-2}\binom{m+i-1}{i}\tau'^{\,i}(1+s+\tau')^{-(m+i)}$ with $s+\tau'=\tau$, which gives \eqref{M-eq:dwellpd} for $\delta>1$. At $\delta=1$ its continuous limit is \eqref{M-eq:dwellpfa}.

\subsection{Self-normalized whitened Doppler}
With a known on-grid Doppler, the whitened dwell vector splits into its DFT coefficient on the target bin, $|\,\cdot\,|^2\sim(1+NS_w|d|^2)E$, and the orthogonal remainder $\Gamma_{N-1}$. The statistic $X/(X+\Gamma_{N-1})$ exceeds $t$ iff $X(1-t)>t\Gamma_{N-1}$, which has probability $\{1+t/((1-t)(1+NS_w|d|^2))\}^{-(N-1)}$; at $S_w=0$ this is $(1-t)^{N-1}$. For the maximum over the $N$ bins, let $I_j$ be independent exponentials with means $\mu_j$ and $J$ a set of bins with $|J|t<1$. Substituting $x_j=tT+y_j$ for $j\in J$, where $T=\sum_iI_i$, the map has Jacobian $1/(1-|J|t)$ and turns the exponent into a sum of independent exponentials, so
\[
\PP\{I_j>tT\ \forall j\in J\}=\frac{1}{1-|J|t}\prod_{i=1}^N\frac{1}{1+\theta_J\mu_i},\qquad\theta_J=\frac{t\sum_{j\in J}\mu_j^{-1}}{1-|J|t},
\]
and zero if $|J|t\ge1$. With all $\mu_i=1$ this is $(1-|J|t)^{N-1}$, and inclusion–exclusion gives Fisher's law $\sum_r(-1)^{r-1}\binom Nr(1-rt)^{N-1}$ (R.~A. Fisher, Proc.\ R.\ Soc.\ Lond.\ A 125, 1929). With one bin of mean $1+NS_w|d|^2$, inclusion–exclusion over sets that do or do not contain it gives $P_{\rm d}$ (\texttt{laws.pd\_snd\_max}).

\subsection{Theorem (certified integration)}
\emph{Domination.} Let $\mathcal J_H$ and $\mathcal J_D$ be the corrupted history and dwell innovations. The $k_H$-th smallest of the history powers is a nondecreasing function of each power, so over all values of the corrupted entries it is smallest when they are zero; hence $s_{k_H}^2\ge\underline s^2$, the value computed with those entries set to zero. If $\underline s^2=0$, define the bound as $W=K\kappa$ (binary: $K$); the bound then holds trivially. For an actual zero scale, choose ratios so each clipped term remains in $[0,\kappa]$ (binary: $[0,1]$); the saturated bound still dominates. Otherwise every uncorrupted dwell term is a clean power divided by $s_{k_H}^2\ge\underline s^2$, and every clipped term is at most $\kappa$. Therefore $T\le W$ for all interference values. For binary integration, an uncorrupted exceedance of the actual statistic is an exceedance at $\underline s^2$, and a corrupted term counts at most one. Adding elements to $\mathcal J$ can only lower $\underline s^2$ and replace terms by their maximum, so $W$ is nondecreasing in $\mathcal J$.
\emph{Calibration.} Draw the calibration masks i.i.d.\ from $\mathcal D$, independently of all clean episodes and of the test hit set. Let the test hit set $\mathcal J$ be independent of all clean episodes and couple it, independently of these episodes and the calibration masks, to $\tilde{\mathcal J}\sim\mathcal D$ with $\mathcal J\subseteq\tilde{\mathcal J}$ (a coupling that exists when $\mathcal D$ dominates the law of $\mathcal J$ in the inclusion order, by Strassen's theorem; a bound on the hit rate alone does not suffice). Then $T\le W(X,\mathcal J)\le W(X,\tilde{\mathcal J})$. The pairs $(X_i,\mathcal J_i)$ and $(X,\tilde{\mathcal J})$ are exchangeable, so the split-conformal rank argument gives $\PP\{W(X,\tilde{\mathcal J})>\widehat q\}\le\alpha$. No property of the interference values is used, so the bound holds for adversarial amplitudes chosen after seeing the clutter, as long as the hit positions are independent of the clutter.
\emph{Trimmed sums.} If a corrupted dwell term can take any value, the trimmed sum over the $K-t$ smallest terms is unbounded once more than $t$ dwell terms are corrupted, which a Bernoulli hit model allows with positive probability. If the history lower bound is positive almost surely, this is exactly the source of infinite dwell terms. The population quantile is infinite when their probability exceeds $\alpha$; a finite-sample threshold is infinite only when enough sampled bounds are infinite. History-scale collapse can also make the trimmed bound infinite.

\subsection{Conformal calibration size and loss}
For continuous i.i.d. scores with $k_\alpha\le n$, the $k_\alpha$-th order statistic of $n$ calibration scores, mapped through the score distribution function, is Beta$(k_\alpha,n+1-k_\alpha)$; the conditional false-alarm probability is one minus it, hence Beta$(n+1-k_\alpha,k_\alpha)$. Because $k_\alpha$ is a ceiling, $\PP\{P_{\rm fa}>2\alpha\}$ is not monotone in $n$; For $\alpha=10^{-2},10^{-3},10^{-4}$, an exhaustive scan gives the smallest $n$ and the $n$ beyond which the condition always holds. The finite scan is completed by a tail bound: writing $L=\lfloor(n+1)\alpha\rfloor$, $\PP\{P_{\rm fa}>2\alpha\}=\PP\{\operatorname{Bin}(n,2\alpha)\le L-1\}$. Since $L-1\le n\alpha$, Chernoff gives an upper bound $e^{-n\alpha/4}<0.05$ for $n\ge12/\alpha$. Thus only the ranks below $12/\alpha$ require enumeration. With ties the rank argument still gives $\PP\{T>\widehat q\}\le\alpha$, but the Beta law no longer applies. For $P_{\rm d}=P_{\rm fa}^{\gamma}$, $\mathbb EU^\gamma=B(n+1-k_\alpha+\gamma,k_\alpha)/B(n+1-k_\alpha,k_\alpha)$, with $B$ the Beta function.
